\section{Introduction}\label{sec:introduction}

The Siegel moduli variety \(\Acal_2\) parametrizes principally polarized abelian surfaces and has dimension three. An indefinite quaternion division algebra can occur as the full rational endomorphism algebra of such a surface. The corresponding non-CM points lie on special curves. The Zilber--Pink conjecture predicts that a Hodge-generic algebraic curve meets the union of these special curves in only finitely many points. The quaternion algebra, its discriminant, and the endomorphism order are allowed to vary in this prediction.

For a geometric point \(s\) of \(\Acal_2\), write \((A_s,\lambda_s)\) for the represented polarized surface and put
\[
 \End^0(A_s)=\End_{\Qbar}(A_s)\otimes_{\Z}\Q.
\]
A curve is \emph{Hodge generic} if it is not contained in a proper special subvariety, including a Hecke translate of one.

\begin{theorem}\label{thm:main}
Let \(C\subset\Acal_{2,\Qbar}\) be a reduced irreducible closed Hodge-generic algebraic curve. Then
\[
 \Sigma_{\mathrm{QM}}(C)=
 \bigl\{s\in C(\Qbar):
   \End^0(A_s)\text{ is an indefinite quaternion division algebra over }\Q
 \bigr\}
\]
is finite.
\end{theorem}

Thus \Cref{thm:main} resolves affirmatively the quaternionic-division component of the Zilber--Pink conjecture for curves in \(\Acal_2\). Its condition is equality of the full geometric endomorphism algebra. In particular, it excludes the split algebra \(M_2(\Q)\) and does not replace the endomorphism condition by the existence of a quaternion embedding into a larger algebra. No condition is imposed on the boundary of \(C\), and the points are counted across all quaternionic special curves at once.

Throughout the paper, a \emph{QM point} means a point whose full geometric rational endomorphism algebra is an indefinite quaternion division algebra over \(\Q\).

\subsection{Earlier results}

The general unlikely-intersection principle asks when an intersection
with a special subvariety has larger dimension than the ambient
dimension count predicts.  Zilber formulated this principle in the
semiabelian setting \cite{Zilber2002}; Pink formulated its Shimura
counterpart \cite[Conjecture~1.3]{Pink2005}.  For a Hodge-generic curve
in a threefold, intersections with special curves have negative
expected dimension.  Pink's nondensity prediction therefore requires
finiteness even after all such special curves are allowed to vary.
The locus in \Cref{thm:main} is the part arising from division
quaternionic multiplication.

Daw and Orr proved the conclusion when the closure of the curve meets the zero-dimensional boundary stratum of the Baily--Borel compactification \cite[Theorem~1.4]{DawOrr2022}. Their reduction for an arbitrary curve assumes a large-Galois-orbit estimate \cite[Theorem~1.3 and Conjecture~6.2]{DawOrr2022}. Their later work proves finiteness against all special curves under that same boundary hypothesis \cite[Corollary~1.2]{DawOrrMultiplicative2025}. Papas obtained further finiteness results under degeneration in higher-dimensional Siegel varieties, with additional conditions on endomorphism algebras and their behavior under reduction; in dimension two, his boundary hypothesis still requires the zero-dimensional stratum \cite[Theorem~1.4]{Papas2024}.

Papas's work with \(v\)-adic values of \(G\)-functions studies neighborhoods of an interior quaternionic point. It gives a height estimate involving the number of supersingular places at which the varying point is close to the chosen center, and finiteness when that number is bounded \cite[Theorem~1.7 and Corollary~1.9]{PapasIII2025}. The horizontal good-reduction comparisons in \cite[Theorem~3.4]{PapasI2025} are also useful here. We prove a height estimate independent of the number of supersingular proximity places.

\subsection{The height argument}

Our height argument belongs to the arithmetic auxiliary-function
tradition of Bombieri and Andr\'e \cite{Bombieri1981,Andre1989}.
Interior period matrices and their $p$-adic realizations enter
Andr\'e's motivic treatment \cite{Andre1995} and Papas's work cited
above.  We prove the interpolation statement needed here directly
for fixed algebraic differential systems at an ordinary point.
In this terminology, ``ordinary'' means that the differential
system has no pole at the center; it imposes no ordinary-reduction
condition on the abelian surface.

To rule out infinitely many targets, pass to a fixed smooth fine-level cover and choose an interior quaternionic point \(t_0\). The substantive estimate is
\begin{equation}\label{eq:intro-height}
 h(t)\le c_1[K(t):K]^{c_2}
 \qquad\bigl(t\text{ a division-quaternionic point}\bigr),
\end{equation}
where \(K\) and the ample logarithmic height \(h\) are fixed. The constants may depend on the curve. This estimate is proved in \Cref{thm:height}; it is not an assumed Galois-orbit bound.

The equations come from a rank-five quadratic variation.  In the first
cohomology of each fiber, take the trace-zero operators that are
self-adjoint for the alternating polarization, with pairing
\(\pair{T}{U}=\tfrac14\Tr(TU)\).  At a QM point, the actual
Rosati-symmetric endomorphisms in this space span a positive definite
two-plane \(E_t\).  We compare a fixed pair of independent integral endomorphisms at the
center with a small integral spanning pair at the target.  Write \(B\) and
\(P\) for their de Rham coordinates, and \(Y(a)\) for horizontal
transport from the target back to the center, where \(a=x(t)\) is a
local algebraic parameter with \(x(t_0)=0\).

At places where the target is close to the center, these coordinates
satisfy a matrix of trace relations
\[
 \pair{B}{Y(a)P}=T_v\in\mathrm M_2(\Q).
\]
The four entries pair the two centered vectors with the two transported
target vectors.  At finite places the relations come from endomorphisms
of their common good reduction.  Although the rational matrix \(T_v\)
can vary with the prime, a lift of residue Frobenius fixes it.  Comparing
the relation with its conjugate therefore eliminates \(T_v\).  The two
arguments are \(a\) and \(b=aR\), where \(b\) is a conjugate of \(a\)
and \(R=b/a\) is included among the algebraic coordinates.  The number of
conjugate tuples is polynomial in the field degree, so one common system retains
a substantial sum of local smallness.  The finitely many exceptional
primes and the archimedean places are grouped directly by their trace
matrices.

The interpolation theorem in \Cref{thm:interpolation} turns this local
information into \eqref{eq:intro-height}.  Its inputs are common analytic
equations, controlled coefficient heights, sufficiently small values of
\(a\), and nonidentity on every algebraic branch used in its proof.
The auxiliary graph is obtained by adjoining the quotient of a finite
Taylor truncation of a common relation by a power of \(a\).  On this
graph, the proof constructs a low-degree polynomial.  The product
formula forces that polynomial to vanish at the target, while the
construction ensures that it does not vanish on the whole graph.
Intersection and projection then reduce the dimension of an algebraic
variety containing the target.  Repeating this descent proves the height
bound.  An ordinary-point differential estimate makes the required orders
polynomial in the algebraic degrees.

It remains to prove the nonidentity input.  Full monodromy handles the
single-argument relations and pairs of arguments that vary independently.
The exceptions are pairs lying on an isogeny correspondence dominating
both fixed curves.  There are only finitely many such correspondences,
up to integral changes of lift.  We choose one auxiliary prime \(\ell\)
and a fixed level cover on which each target admits a marking with the
following property: modulo \(\ell\), its plane and the centered plane
span a degenerate four-space.  A functional identity in the exceptional
case would force a Frobenius quasi-isogeny to act by one common sign on
both planes.  The prescribed incidence rules this out, using the
polarized multiplier and the integral trace form.  This argument also
covers supersingular reductions and target planes that are themselves
degenerate modulo \(\ell\).

Once \eqref{eq:intro-height} is known, endomorphism estimates and
one-variable Hodge-norm bounds give polynomially bounded integral
coordinates for two Hodge vectors at every point of a large Galois orbit.
The periods lie in the three-dimensional quadric of isotropic lines in
the rank-five space.  The period at a QM point is orthogonal to its
plane \(E_t\), and the orthogonal complement of each positive two-plane
cuts this quadric in a conic.  The path form of o-minimal counting would produce a one-parameter
semialgebraic family of such planes whose conics contain a nonconstant
period arc.  Their union has algebraic closure of dimension at most two.
Full monodromy, however, makes the period image Zariski dense in the
three-dimensional quadric, giving the contradiction.

This last step follows the arithmetic point-counting strategy of
Pila and Zannier \cite{PilaZannier2008}.  The rational-point theorem of
Pila and Wilkie \cite{PilaWilkie2006}, refined into definable blocks
by Pila \cite[Theorem~3.6]{Pila2011}, underlies the projected
path-counting result of Habegger and Pila
\cite[Corollary~7.2]{HabeggerPila2016} used here.  The numerical
bounds for the Hodge vectors come from the endomorphism estimates of
Masser and W\"ustholz \cite{MW1994} and Schmid's degeneration and norm
theorems \cite{Schmid1973}.  In the rank-five variation the final
geometric contradiction is a direct dimension calculation followed
by the monodromy theorem of Andr\'e \cite{Andre1992}.

\subsection{Conventions and organization}

All varieties and algebraic points are over \(\Qbar\) unless a field is specified. Algebraic heights are absolute logarithmic heights. Local absolute values extend the usual, nonsquared absolute values on \(\Q\), and local sums carry the normalized weights. The rational coordinate height used in point counting is multiplicative and will be denoted \(H_{\mathrm{rat}}\).

Constants and polynomial exponents may depend on fixed curves, covers, centers, frames, and connections. They never depend on the varying quaternion algebra or target point. A polynomial bound always refers to a fixed number of variables and fixed connection ranks.

\Cref{sec:geometry} establishes the geometric and monodromy input. The interpolation theorem is proved in \Cref{sec:interpolation}. \Cref{sec:arithmetic} constructs the arithmetic data and common equations. \Cref{sec:height} proves their nonidentity and the height estimate. \Cref{sec:counting} finishes the proof and descends to the coarse curve.
